% Terminal note of the homily: everything the speech itself must not carry ---
% the audience, the spoken word count and the pace it implies, the relation to
% the three readings of the expansive study, the route by which each quoted
% text reaches the page, and the exact loci --- followed by the References.
% Nothing here is read aloud.

\section*{Note on Sources and Delivery}
\runninghead{Sources and Delivery}

{\footnotesize

\textbf{Audience and occasion.} An adult parish assembly, at the Mass
\latin{Dominica decima nona post Pentecosten}, II classis, of the 1962
\work{Missale Romanum}, editio typica, pp.~411--412, nos.~1679--1688, for the
occurrence of 4~October 2026.

\textbf{Length and pace.} The spoken body runs to 1,453 words, about 11.2 to
12.1 minutes at an unhurried 120 to 130 words a minute by arithmetic alone; no
one has delivered or timed it. It was read through in full, silently, for sense
and ease of speech.

\textbf{Relation to the reviewed readings.} The homily joins the first and third
of the expansive study's three readings of the whole Mass: \textit{The Garment
the King Looks For} (Gregory's and Augustine's charity, Jerome's new man who is
Christ, the Epistle's four acts, the bound and the lifted hands, and the
Communion's wish as a prayer for grace) and \textit{The Feast That Now Is}
(Augustine's two feasts, his table that is not the garment, and his evening
sacrifice of the Passion). From the second, \textit{The Invitation Refused and
Carried to the Highways}, it takes only the story of the call and Chrysostom's
two sentences; it does not say who the first invited or the highways were. It
speaks two disagreements the study keeps: Augustine against Hilary on the
evening sacrifice, and Hilary's baptismal gift beside Gregory's and Augustine's
charity. Jerome joins the new man to the garment himself. The other joins are
the homily's own and credited to no Father: the Collect beside the excuses, the
charity the Bridegroom wore on the Cross, the Communion's wish against the one
Augustine forbids, and the Introit's promise beside the asking. So are the
sentences on the altar, the Church's faith in the preacher's own voice; no
Father is said to call the wedding feast the Eucharist. The Alleluia, Offertory
and Secret are left to the studies.

\textbf{Texts quoted, and their route.} Scripture is the Douay--Rheims,
Challoner's revision, at the canonical verses. The Introit's antiphon, the
Collect and the Postcommunion are quoted in the anonymous English printed by
Eugene Cummiskey at Philadelphia in 1861, a historical study aid and not an
approved liturgical translation of the 1962 Missal; the Postcommunion's
``healing work'' names the Latin \latin{medicinalis operatio}, and the 1861
``these thy mysteries'' is not quoted. Augustine and Chrysostom are quoted in
the English of the Nicene and Post-Nicene Fathers. Gregory, Jerome and Hilary
are reported in English paraphrase of their Latin, and their words are not
quoted.

\textbf{Exact loci.} Introit, no.~1679; Collect, no.~1680; Epistle,
Eph 4:24--28, no.~1681; Gradual, Ps 140:2, no.~1682; Gospel, Mt 22:3, 5, 10,
12, no.~1684; Communion, Ps 118:4--5, no.~1687; Postcommunion, no.~1688; with
Mt 25:35 and Ps 138:16. Chrysostom, \work{Homilies on Matthew} 69. Gregory,
\work{Homiliae in Evangelia} 38, 9 (not baptism or faith, but charity), 11
(hatred and envy), 12 (friend by faith; Ps 138:16), 13 (the bound hands) and 14
(called and chosen). Augustine, \work{Sermo} 90, 1, 4, 5 and 9, and \work{Sermo}
95, 7; \work{Enarrationes in Psalmos} 118, Aleph 5, and 140, 3. Jerome on
Mt 22:11--12 (PL 26, col.~160) and on Eph 4:24 (col.~508). Hilary on Mt 22:11
(PL 9, col.~1044) and on Ps 140:2 (cols.~825--826).

\par}

\Needspace{14\baselineskip}
\section*{References}
\runninghead{References}

{\footnotesize
\begin{itemize}\setlength{\itemsep}{0.15em}\setlength{\parskip}{0pt}
\item \work{Missale Romanum ex decreto Sacrosancti Concilii Tridentini restitutum, Summorum Pontificum cura recognitum}, editio typica (Typis Polyglottis Vaticanis, 1962), CMAA facsimile, pp.~411--412, nos.~1679--1688.
\item \work{The Roman Missal, Translated into the English Language for the Use of the Laity} (Philadelphia: Eugene Cummiskey, 1861), pp.~445--448 (XIX Sunday after Pentecost, Introit, Collect, Postcommunion).
\item The Holy Bible, Douay--Rheims version, Challoner's revision: Ps 118:4--5; 138:16; 140:2; Mt 22:3, 5, 10, 12; 25:35; Eph 4:24--28.
\item St John Chrysostom, \work{Homilies on Matthew} 69, trans. G. Prevost, rev. M.~B. Riddle, Nicene and Post-Nicene Fathers, 1st ser., vol.~10.
\item St Gregory the Great, \work{Homiliae in Evangelia} 38, 9 and 11--14, Latin Wikisource transcription of Migne's text; paraphrased, not quoted.
\item St Augustine, \work{Sermones} 90 and 95 (Sermons XL and XLV), Nicene and Post-Nicene Fathers, 1st ser., vol.~6: Sermo 90, 1, 4--5, 9; Sermo 95, 7. \work{Enarrationes in Psalmos} 118, Aleph 5, and 140, 3 (Psalms CXIX and CXLI in that translation), 1st ser., vol.~8.
\item St Jerome, \work{Commentarii in Matthaeum} III and \work{Commentarii in Epistolam ad Ephesios} II, in Migne, PL 26 (Paris, 1845), cols.~160 and 508; paraphrased, not quoted.
\item St Hilary of Poitiers, \work{Commentarius in Matthaeum} 22, 7, and \work{Tractatus super Psalmos}, in Ps.~140, 3--4, in Migne, PL 9 (Paris, 1844), cols.~1044 and 825--826; paraphrased, not quoted.
\end{itemize}
\par}
